\cleardoublepage
\chapter*{Introduction générale}
\addcontentsline{toc}{chapter}{Introduction générale}

Le secteur bancaire marocain constitue le pilier du financement de l’économie nationale et joue un rôle de catalyseur pour la croissance des agents économiques. Dans un environnement macroéconomique marqué par des fluctuations conjoncturelles et un renforcement strict des exigences réglementaires imposées par Bank Al-Maghrib, les établissements de crédit sont contraints d'optimiser en permanence la gestion de leurs expositions. L'intermédiation financière, cœur de métier des banques, implique structurellement une prise de risque qui nécessite des dispositifs d'évaluation de plus en plus sophistiqués pour préserver la rentabilité et les fonds propres.

Au sein de ce secteur, la filière du financement des entreprises (Corporate Banking) représente à la fois la principale source de revenus et le vecteur de risque le plus significatif. L'octroi de facilités de caisse, de lignes d'escompte ou de cautions administratives et diverses expose mécaniquement l'institution au risque de contrepartie. Face à cette dynamique sectorielle, la maîtrise du risque de crédit est devenue un avantage concurrentiel décisif, imposant le passage de méthodes d'évaluation empiriques à des systèmes de notation interne (rating) quantitatifs et rigoureux.

C’est dans ce cadre sectoriel qu'évolue Bank of Africa (Groupe BMCE). Fondée le 1er septembre 1959 sous le nom de Banque Marocaine du Commerce Extérieur (BMCE) pour accompagner le commerce international du Maroc post-indépendance, l'institution a suivi une trajectoire d'expansion continue, jalonnée par sa privatisation en 1995 et son déploiement panafricain à partir de 2008. Le passage à la marque unique "Bank of Africa" en 2019 a marqué la volonté du groupe de consolider son ancrage multinational et d'unifier ses forces sur le continent.

Aujourd'hui, l'établissement se positionne comme un groupe bancaire universel multi-métiers (banque commerciale, banque d'affaires, services financiers spécialisés), s'appuyant sur un capital humain de près de 15300 collaborateurs pour servir 6,6 millions de clients à travers un réseau mondial de 2000 points de vente. Ses orientations stratégiques actuelles privilégient l'intégration de la banque digitale, le soutien aux PME africaines et le renforcement des dispositifs de gestion des risques. Pour répondre efficacement aux besoins de la clientèle "Entreprises", le réseau s'articule autour de structures dédiées. Le présent stage de fin d'année se déroule au sein du Centre d'Affaires de Bank of Africa à Agadir. Cette entité centralise le traitement, l'analyse et le suivi des engagements corporate (crédits à court terme, gestion des effets de commerce et émission de cautions), offrant un ancrage pratique idéal pour étudier les processus d'octroi de crédit.

L'activité de financement des entreprises expose le Centre d'Affaires au risque de défaut. Ce constat soulève une problématique financière majeure, au cœur de la présente étude : L'analyse comparative entre les modèles de rating de risque de crédit. Plusieurs interrogations émergent alors : Comment se manifeste concrètement l'évaluation du risque de crédit lors de l'étude d'un dossier de financement ? Quels sont les fondements et les spécificités des différents modèles de rating (modèles structurels, statistiques ou algorithmiques) actuellement mobilisables par les institutions financières ? Dans quelle mesure une analyse comparative rigoureuse permet-elle d'identifier les modèles les plus performants, précis et adaptés pour prédire le défaut d'une contrepartie dans le contexte économique actuel ?

Pour aborder efficacement ce sujet, nous allons mobiliser un cadre conceptuel structuré autour de trois axes principaux. Le premier porte sur la compréhension globale du risque de crédit dans la théorie de l'intermédiation financière et le cadre réglementaire international, notamment les directives de Bâle sur les approches de notation interne (IRB). Le second axe concerne la nature même de la notation financière (rating) et les variables explicatives de la solvabilité d'une entreprise. Le troisième s’intéresse aux différents modèles quantitatifs d'évaluation du risque (modèles discriminatoires comme le Z-score d'Altman, modèles logistiques, ou approches probabilistes). Ce cadre théorique permet d'étayer l'analyse de la situation et d'établir un référentiel solide pour la comparaison des modèles.

Afin de répondre à ces questions, une méthodologie combinant approche qualitative et analyse quantitative a été adoptée. Sur le plan qualitatif, l'immersion au sein des services opérationnels du Centre d'Affaires et des entretiens avec les responsables ont permis de mieux cerner les pratiques de sélection des risques en place et les contraintes opérationnelles perçues. Sur le plan quantitatif, une analyse basée sur des données financières d'entreprises permettra de simuler l'application de différents modèles de rating. Cette double approche permet de croiser les processus internes avec la théorie mathématique, renforçant ainsi la pertinence de l'analyse comparative.

Ce stage s’inscrit dans le cadre du cursus académique de l’École Nationale de Commerce et de Gestion (ENCG), dont la vocation est de former des cadres capables d’articuler théorie et pratique. Pour l'entreprise, l'objectif est d'obtenir une étude comparative documentée sur les outils de modélisation, pouvant offrir des perspectives d'amélioration des processus de décision d'octroi de crédit. Pour moi, il s'agit de mobiliser les connaissances en finance quantitative, de développer des compétences techniques et analytiques sur la modélisation des risques, et d'affiner son projet professionnel. Pour l'École, ce travail valide la capacité de ses étudiants à traiter des problématiques financières complexes inhérentes au marché du travail.

Le présent rapport s’articule en trois chapitres. Le premier expose le cadre conceptuel de l’étude, dédié à la théorie du risque de crédit et aux modèles de rating. Le deuxième est consacré à la présentation de Bank Of Africa et du Centre d'Affaires d'Agadir, ainsi qu'à la méthodologie de recherche adoptée. Le troisième présente les résultats de l'analyse comparative des modèles, les recommandations qui en découlent, ainsi qu’un bilan personnel et professionnel du stage.